\par\noindent\textbf{(1)} 电子线性方程 \(m(\partial_t+\nu)\boldsymbol v=-e(\boldsymbol E+\boldsymbol v\times B_0\hat{\boldsymbol z})\)。以 \(\boldsymbol E\propto e^{-i\omega t}\) 并取圆偏振 \(v_\pm=v_x\pm iv_y\)，得到 \([\nu-i(\omega\mp\omega_c)]v_\pm=-(e/m)E_\pm\)，其中 \(\omega_c=eB_0/m\)。
\par\noindent\textbf{(2)} 因为 \(\boldsymbol J=-ne\boldsymbol v\)，复电导为 \(\sigma_\pm(\omega)=ne^2/[m(\nu-i(\omega\mp\omega_c))]\)。由 \(\boldsymbol D=\varepsilon_0\boldsymbol E+i\boldsymbol J/\omega\) 得 \[\varepsilon_\pm/\varepsilon_0=1+i\sigma_\pm/(\varepsilon_0\omega)=1-\omega_p^2/[\omega(\omega\mp\omega_c+i\nu)]\]，\(\omega_p^2=ne^2/(\varepsilon_0m)\)。
\par\noindent\textbf{(3)} 时间平均耗散功率密度为 \(\langle p_\pm\rangle=\tfrac12\operatorname{Re}\sigma_\pm|E_\pm|^2\)，其中 \(\operatorname{Re}\sigma_\pm=\frac{ne^2\nu/m}{\nu^2+(\omega\mp\omega_c)^2}\ge0\)。适当手性的共振峰位于 \(\omega\approx\omega_c\)，半高全宽 \(2\nu\)（角频率）。
\par\noindent\textbf{(4)} 碰撞趋零时Lorentz吸收峰宽趋零，其积分谱权重有限，远离共振处耗散消失；反向 \(B_0\) 则交换两种圆偏振的峰。近回旋共振还可能有动理学波粒相互作用，不能以此冷单流体谱描述所有高温等离子体。
